\clearpage
\section{Notation}\label{app:notation}

This section collects the asymptotic conventions and the principal symbols
used in the main text.  The tables follow the order of the sections
where the corresponding symbols that are mainly introduced.
Temporary variables used only within individual proofs are omitted.
The domain or regime column records the scope of each definition.

\subsection{Asymptotic notation}\label{app:asymptotic-notation}

Subscripts on $O$, $\Omega$, and $\Theta$ specify the dependence
of implicit constants; unsubscripted $O$-terms have absolute constants.
Set subscripts, such as $K$ and $U$ in $O_{H,K}$ and $O_{H,U}$,
allow uniform dependence over the parameters in that set.
For $o$-terms, subscripts indicate fixed parameters; convergence is taken
in the stated regime; no uniformity in these parameters is asserted.
The notation $f\ll g$ means $f/g\to0$, and $f\gg g$ means $g\ll f$.

\subsection{Symbols}\label{app:symbols}

All integrals without an indicated measure use Lebesgue measure, usually on $[0,1]$ or $[0,1]^2$.  The notation $\|\cdot\|_k$
denotes the corresponding $L^k$-norm.  We use natural logarithms and
$0\log0=0$.

\newenvironment{notationtable}{%
  \small
  \setlength{\tabcolsep}{4pt}%
  \renewcommand{\arraystretch}{1.18}%
  \setlength{\LTpre}{4pt}%
  \setlength{\LTpost}{6pt}%
  \begin{longtable}{@{}>{\raggedright\arraybackslash}p{0.17\textwidth}
    >{\raggedright\arraybackslash}p{0.21\textwidth}
    >{\raggedright\arraybackslash}p{\dimexpr0.62\textwidth-4\tabcolsep\relax}@{}}
  \toprule
  Symbol & Domain or regime & Definition and role \\
  \midrule
  \endfirsthead
  \toprule
  Symbol & Domain or regime & Definition and role \\
  \midrule
  \endhead
  \bottomrule
  \endfoot
}{\end{longtable}}

\subsubsection*{\Cref{sec:introduction}: Basic objects and the variational problem}

\begin{notationtable}
$H,d,m,v$ & Main regular-graph results &
$H$ is a finite simple $d$-regular graph, $d\ge2$;
$m=|E(H)|$ and $v=|V(H)|=2m/d$.  The symbol $v$ is introduced in
\Cref{sec:singular-endpoint}. \\
$G(n,p),X_H$ & $n\in\mathbb N$, $0<p<1$ &
The Erd\H{o}s--R\'enyi random graph and its number of injective
homomorphisms from $H$, respectively. \\
$(n)_k$ & Integers $0\le k\le n$ &
$n(n-1)\cdots(n-k+1)$, the falling factorial. \\
$p,r$ & Dense regime: $0<p<r<1$ &
The background edge probability and target density, fixed as
$n\to\infty$.  The target $H$-density is $r^m$. \\
$t(H,G)$,
$\mathcal E_{H,r}$ & A finite graph $G$ &
$t(H,G)=\hom(H,G)/|V(G)|^{|V(H)|}$;
$\mathcal E_{H,r}=\{t(H,G(n,p))\ge r^m\}$ is the upper-tail event. \\
$W,\W_0$ & $W:[0,1]^2\to[0,1]$ &
A symmetric measurable function (a graphon), and the space of all
graphons.  A scalar $r$ also denotes the constant graphon $W\equiv r$. \\
$e(W),t(H,W)$ & $W\in\W_0$ &
$e(W)=\int W=t(K_2,W)$; $t(H,W)$ is the integral of
$\prod_{ij\in E(H)}W(x_i,x_j)$ over the vertex variables. \\
$J_p(z),I_p(W)$ & $0<p<1$, $0\le z\le1$ &
$J_p(z)=z\log(z/p)+(1-z)\log((1-z)/(1-p))$;
$I_p(W)=\int J_p(W)$ is the graphon relative entropy. \\
$\Phi_H(p,r)$,
$\W_H(p,r)$ & $0<p<r<1$ &
The minimum of $I_p(W)$ subject to $t(H,W)\ge r^m$, and its set of
minimizers; see \eqref{eq:graphon-variational-problem}. \\
$\lambda(p,r)$ & $0<p,r<1$ &
$\log((1-p)/p)-\log((1-\pc(r))/\pc(r))$;
the log-odds displacement from the phase boundary, positive when
$p<\pc(r)$. \\
\end{notationtable}

\subsubsection*{\Cref{sec:preliminaries}: Graphon topology and entropy maximization}

\begin{notationtable}
$\|F\|_\square$ & Integrable kernels on $[0,1]^2$ &
$\sup_{S,T}|\int_{S\times T}F|$, with measurable $S,T\subseteq[0,1]$;
the cut norm. \\
$W^\sigma$,
$\delta_\square(U,W)$ & $U,W\in\W_0$ &
$W^\sigma(x,y)=W(\sigma x,\sigma y)$ and
$\delta_\square(U,W)=\inf_\sigma\|U-W^\sigma\|_\square$.
Here $\sigma$ is a measure-preserving a.e.\ bijection. \\
$\widetilde\W_0,W_G$ & Graphons; finite graphs $G$ &
The graphon space modulo zero cut distance, and the step
graphon associated with $G$.  Uniqueness is understood in
$\widetilde\W_0$; see \Cref{sec:graphons}. \\
$q_{11},q_{12},q_{22}$;
$c$ & Bipodal graphons;
$q_{ij},c\in[0,1]$ &
The two within-block densities, cross-block density, and first-block
size: $q_{11}$ belongs to the block of size $c$ and $q_{22}$ to its
complement.  In the KRR--S parametrization the smaller block is first. \\
$S_0(z),s(W)$ & $z\in[0,1]$, $W\in\W_0$ &
$S_0(z)=-\tfrac12[z\log z+(1-z)\log(1-z)]$;
$s(W)=\int S_0(W)$. \\
$\eps,\tau,\vartheta$ & Fixed-density slices &
$\eps=e(W)$, $\tau=t(H,W)$, and $\vartheta=\tau-\eps^m$. \\
$\psi_d(\eps,z)$ & $\eps,z\in(0,1)$ &
$\displaystyle\frac{2[S_0(z)-S_0(\eps)-S_0'(\eps)(z-\eps)]}
{z^d-\eps^d-d\eps^{d-1}(z-\eps)}$, with the diagonal value filled in
continuously; the scalar quotient selecting the limiting cross density. \\
$\zeta_d(\eps)$ & $\eps\in(0,1)$ &
The unique maximizer of $z\mapsto\psi_d(\eps,z)$.
For $\eps\ne r_*$, it is the limiting KRR--S cross density and
satisfies $\zeta_d(\eps)\ne\eps$. \\
\end{notationtable}

\subsubsection*{\Cref{sec:lz-boundary}: The Lubetzky--Zhao boundary and its two contacts}

\begin{notationtable}
$r_*,p_*$ & Fixed by $d$ &
$r_*=(d-1)/d$ and
$p_*=(d-1)/(d-1+\exp(d/(d-1)))$;
the exceptional target density and boundary probability. \\
$\varphi_{p,d}(x)$ & $p\in(0,1)$, $x\in[0,1]$ &
$J_p(x^{1/d})$; its convex minorant determines the phase boundary. \\
$h_{p,d}(z)$ & $z\in(0,1)$ &
$zJ_p''(z)-(d-1)J_p'(z)$; its sign is the sign of
$\varphi_{p,d}''(z^d)$. \\
$x_a(p),x_b(p)$ & $0<p<p_*$ &
The two contact points of the affine part of the convex minorant.
Their density coordinates are $u_a=x_a^{1/d}$ and $u_b=x_b^{1/d}$;
see \Cref{lem:contact-points}. \\
$\pc(r),\sm(r)$ & $r\in(0,1)$ &
The phase-boundary probability and the other contact density.
They are analytic for $r\ne r_*$ and extend continuously by
$\pc(r_*)=p_*$, $\sm(r_*)=r_*$.  For $r\ne r_*$,
$\sm(r)=\zeta_d(r)$; see \Cref{rmk:bipodal-parameter-expansions}. \\
$\ell_r(x)$ & $x\in[0,1]$ &
$J_{\pc(r)}(r)+[J_{\pc(r)}'(r)/(dr^{d-1})](x-r^d)$;
the supporting tangent, touching at $r^d$ and $\sm(r)^d$ when
$r\ne r_*$.  It extends to $r=r_*$ in \Cref{sec:constant-graphon-comparison}. \\
\end{notationtable}

\subsubsection*{\Cref{sec:nonexceptional-endpoint}: Nonexceptional optimizers}

\noindent\emph{Reduction at fixed densities
(\Cref{sec:local-reduction}).}
\par\nopagebreak[4]

\begin{notationtable}
$I$ & Open interval containing $r_0$, with compact closure in
$(0,1)\setminus\{r_*\}$ &
A sufficiently small neighborhood of the fixed density $r_0$ on which
the nonexceptional estimates are uniform; it may be shrunk as needed. \\
$S_H(\eps,\tau)$ & $(\eps,\tau)\in[0,1]^2$ &
$\sup\{s(W):e(W)=\eps,\ t(H,W)=\tau\}$;
the entropy envelope, equal to $-\infty$ for an empty slice. \\
$I_{p,r}(\eps)$ & $p,r\in(0,1)$,
$\eps\in[0,1]$ &
$-2S_H(\eps,r^m)-\log(1-p)+\eps\log((1-p)/p)$;
the relative-entropy infimum at densities $(\eps,r^m)$, equal to
$+\infty$ for an empty slice. \\
$B_{\eps,r}$ & $r\in I$, $\eps<r$ sufficiently close to $r$ &
The unique bipodal entropy maximizer at densities $(\eps,r^m)$;
$I_p(B_{\eps,r})=I_{p,r}(\eps)$.
See \Cref{lem:fixed-density-bipodality}. \\
\end{notationtable}

\noindent\emph{The boundary cost and its curvature
(\Cref{sec:nonexceptional-quadratic-growth}).}
\par\nopagebreak[4]

\begin{notationtable}
$\delta$ & $r\in I$ &
$\delta=r-\eps$ is the edge-density deficit. \\
$g_r(z)$ & $r\in I$, $z\in[0,1]$ &
$J_{\pc(r)}(z)-\ell_r(z^d)$;
the nonnegative supporting cost gap, vanishing exactly at
$r$ and $\sm(r)$; see \eqref{eq:supporting-cost-gap}. \\
$D_d(r,z)$ & $r,z\in(0,1)$, $z\ne r$ &
$\displaystyle\frac{z^d-r^d}{dr^{d-1}}-(z-r)>0$;
the coefficient relating smaller-block size to edge-density deficit, used in
the comparison construction and \eqref{eq:edge-deficit-coefficient}. \\
$G_r(\delta)$,
$G(r,\delta)$ & Small $\delta>0$;
respectively, a neighborhood of $I\times\{0\}$ &
$G_r(\delta)=I_{\pc(r),r}(r-\delta)-J_{\pc(r)}(r)$
is the boundary excess cost.  $G$ is its two-sided analytic extension;
derivatives at $\delta=0$ refer to $G$. \\
$A_H(r)$ & $r\in(0,1)\setminus\{r_*\}$ &
$\partial_\delta^2G(r,0)>0$;
the analytic coefficient of the quadratic boundary cost, introduced
in \Cref{thm:positive-second-variation}. \\
\end{notationtable}

\noindent\emph{The nonexceptional optimizer
(\Cref{thm:nonexceptional-endpoint}; proof in
\Cref{sec:nonexceptional-proof}).}
\par\nopagebreak[4]

\begin{notationtable}
$\delta_*(p,r)$ & Near $p=\pc(r)$, $r\ne r_*$ &
The analytic solution of
$\partial_\delta G(r,\delta_*)=\lambda(p,r)$.
On the symmetry-breaking side it equals $r-e(W_{p,r})$ and is positive. \\
$W_{p,r}$ & Near the nonexceptional boundary, $p<\pc(r)$ &
The unique optimizer, up to relabeling;
$W_{p,r}=B_{r-\delta_*(p,r),r}$ up to relabeling. \\
$A_{p,r},c(p,r)$;
$q_{ij}(p,r)$ & The same regime &
$A_{p,r}=[0,c(p,r)]$ is the smaller block, with
$0<c(p,r)<1/2$.  The three densities and $c$ are the KRR--S parameters
evaluated at $\eps=r-\delta_*$,
$\vartheta=r^m-(r-\delta_*)^m$. \\
\end{notationtable}

\subsubsection*{\Cref{sec:singular-endpoint}: The exceptional family and its comparison}

\noindent\emph{Stationary family and scalar cost
(\Cref{sec:rank-one-stationary-family,sec:constant-graphon-comparison}).}

\begin{notationtable}
$u_*,P_*$ & Fixed by $d$ &
$u_*=\sqrt{r_*}$ and $P_*=(p_*,r_*)$;
the limiting factor value and exceptional boundary point. \\
$\lambda$,
$f\otimes f$ & \Cref{sec:singular-endpoint} &
$\lambda$ is Lebesgue measure, not $\lambda(p,r)$.
$(f\otimes f)(x,y)=f(x)f(y)$ is a rank-one kernel; it need not be a
graphon for a general bounded factor $f$. \\
$q(f)$ & Nonnegative $f\in L^d([0,1])$ &
$\int_0^1f^d$; regularity gives $t(H,f\otimes f)=q(f)^v$.
In the localization and comparison proofs, $q$ abbreviates $q(f)$. \\
$f_{\alpha,s,t}$,
$W_{\alpha,s,t}$ & $\alpha,s,t\in(0,1)$ &
$f_{\alpha,s,t}=s\1_{[0,\alpha]}+t\1_{(\alpha,1]}$;
$W_{\alpha,s,t}=f_{\alpha,s,t}\otimes f_{\alpha,s,t}$. \\
$u,h,s_h,t_h$ & Near $(u,h)=(u_*,0)$ &
$u=(s+t)/2$ and $h=(t-s)/2$.
Along the family, $s_h=u_h-h$ and $t_h=u_h+h$;
for $h>0$ the two factor values satisfy $s_h<t_h$. \\
$\alpha_h,f_h,W_h$ & $|h|<h_0$ &
$f_h=f_{\alpha_h,s_h,t_h}$ and $W_h=f_h\otimes f_h$.
$\alpha_h$ is the mass of the $s_h$-valued block, not necessarily the
smaller block; $\alpha_{-h}=1-\alpha_h$ and $\alpha_0=1/2$. \\
$p_h,q_h,r_h$ & $|h|<h_0$ &
The background probability, factor moment
$q_h=q(f_h)=\alpha_hs_h^d+(1-\alpha_h)t_h^d$, and target
$r_h=q_h^{2/d}$; thus $t(H,W_h)=q_h^v=r_h^m$. \\
$\gamma_h,\mu_h$,
$\gamma_*$ & $|h|<h_0$ &
$\gamma_h=\mu_h m q_h^{v-2}>0$ is the scalar stationarity coefficient;
$\mu_h$ is the graphon constraint multiplier.
$\gamma_*=r_*^{-d}$ is the value of $\gamma_h$ at $h=0$. \\
$F_{p,\gamma}(z)$ & $z\in(0,1)$ &
$J_p'(z)-\gamma z^{d-1}$; stationarity requires
$F_{p,\gamma}(f(x)f(y))=0$ a.e., with
$\gamma=\mu m q(f)^{v-2}$; see \eqref{eq:rank-one-kkt}. \\
$\ell(z),\ell_*$,
$\Lambda_h$ & $z\in(0,1)$ &
$\ell(z)=\log((1-z)/z)$, $\ell_*=\ell(p_*)$,
$\Lambda_h=\ell(p_h)-\ell_*$.
$\Lambda_h$ is the slope of the affine change in
$J_{p_h}-J_{p_*}$, not the displacement $\lambda(p_h,r_h)$. \\
$\mathcal M_h(z)$,
$R(h,z)$ & $z$ near $r_*$ &
$\mathcal M_h(z)=J_{p_h}(z)-(\gamma_h/d)z^d$.
The analytic quotient $R$ is defined by
$\mathcal M_h'(z)=(z-s_h^2)(z-s_ht_h)(z-t_h^2)R(h,z)$;
see \eqref{eq:mh-derivative-factorization}. \\
$\beta_d,\Gamma_d(z)$ & $z\in[0,1]$ &
$\beta_d=\gamma_*/d$ is the tangent slope at the singular endpoint;
$\Gamma_d(z)=J_{p_*}(z)-\ell_{r_*}(z^d)$ is the corresponding supporting
cost gap.  It vanishes to fourth order at $r_*$ and controls
$|z-r_*|^4$; see \eqref{eq:endpoint-supporting-gap}. \\
\end{notationtable}

\noindent\emph{Localization and decomposition
(\Cref{sec:localization-rank-one}).}
\par\nopagebreak[4]

\begin{notationtable}
$d',T_K$ & $d'=d/(d-1)$;
suitable kernels $K$ &
$d'$ is the H\"older conjugate of $d$, and
$(T_Kg)(x)=\int_0^1K(x,y)g(y)\,dy$.
Here $K$ denotes a kernel, not the compact density interval. \\
$\mathcal Q_W,g_0$,
$\theta$ & $g\in L^{d'}([0,1])$, $\|g\|_{d'}\le1$ &
$\mathcal Q_W(g)=\iint W(x,y)g(x)g(y)\,dx\,dy$;
$g_0\ge0$ is a unit-norm maximizer and
$\theta=\mathcal Q_W(g_0)>0$. \\
$f,E$ & Competitors in \Cref{lem:localization-rank-one} &
$f=\theta^{1/2}g_0^{1/(d-1)}$ and $E=W-f\otimes f$.
They satisfy $0\le f\le2$ and $T_E(f^{d-1})=0$;
$\|f-u_*\|_4,\|E\|_2=O_d(h)$. \\
\end{notationtable}

\noindent\emph{Factor distributions and Lagrangian comparison
(\Cref{sec:auxiliary-lagrangian,sec:graphon-comparison}).}
\par\nopagebreak[4]

\begin{notationtable}
$\nu,\nu_h$ & Probability measures on $[0,2]$ &
$\nu=f_\#\lambda$ is the distribution of factor values;
$\nu_h=(f_h)_\#\lambda=\alpha_h\delta_{s_h}+(1-\alpha_h)\delta_{t_h}$,
where $\delta_x$ is the Dirac mass at $x$. \\
$\xi,\sigma$ & Measures on $[0,2]$ &
$\xi$ is a general probability measure and $\sigma=\xi-\nu_h$ is its
signed difference from the reference measure; for a competitor,
$\xi=\nu$. \\
$m_d(\xi),\Delta_h(\xi)$ & Probability measures $\xi$ on $[0,2]$ &
$m_d(\xi)=\int x^d\,d\xi(x)$ and
$\Delta_h(\xi)=m_d(\xi)-q_h$;
in particular $m_d(\nu)=q(f)$ and $m_d(\nu_h)=q_h$. \\
$\widetilde\Gamma_d$,
$\widetilde J_{p_h}$ & Scalar arguments in $[0,4]$ &
The auxiliary continuations in \eqref{eq:entropy-continuation}.
They agree with $\Gamma_d,J_{p_h}$ on $[0,1]$;
above $1$, $\widetilde\Gamma_d(z)=\Gamma_d(1)+(z-1)^2$.
The continued cost is not a relative entropy above $1$. \\
$K_h(x,y)$ & $(x,y)\in[0,2]^2$ &
$\widetilde J_{p_h}(xy)$; the cost kernel in factor-value coordinates,
as opposed to the graphon kernel on the vertex space. \\
$\mathcal J_h(\xi)$,
$\mathcal L_h(\xi)$ & Probability measures $\xi$ on $[0,2]$ &
$\mathcal J_h(\xi)=\iint K_h\,d\xi\,d\xi$;
$\mathcal L_h(\xi)=\mathcal J_h(\xi)-\mu_h\{m_d(\xi)^v-q_h^v\}$.
These are the auxiliary factor functional and Lagrangian. \\
$\Psi_h(x),\kappa_h$ & $x\in[0,2]$ &
$\Psi_h(x)=2\int K_h(x,y)\,d\nu_h(y)
-(2\gamma_hq_h/d)x^d-\kappa_h$, with $\Psi_h(s_h)=0$.
It represents the first variation of $\mathcal L_h$ at $\nu_h$;
see \eqref{eq:first-variation}. \\
$\mathcal N_\rho$,
$\mathcal T_\rho$ & Small fixed $\rho>0$ &
$\mathcal N_\rho=(u_*-\rho,u_*+\rho)$ and
$\mathcal T_\rho=[0,2]\setminus\mathcal N_\rho$;
the central and tail regions in factor-value space. \\
$Q_h(x)$,
$A_{h,\rho}(\xi)$ & $x\in[0,2]$;
probability measures $\xi$ &
$Q_h(x)=(x-s_h)(x-t_h)$ and
$A_{h,\rho}(\xi)=\int_{\mathcal N_\rho}Q_h^2\,d\xi$;
the central deviation from the two reference factor values. \\
$I_h^x,R_h^x$ & Functions of $(x,y)$ near $(u_*,u_*)$;
$h>0$ &
Interpolation in $x^d$ at $s_h^d,t_h^d$, and the remainder
$R_h^xF=F-I_h^xF$; the $y$-operators are analogous.
See \eqref{eq:moment-interpolation}. \\
$\Omega_\rho$ & A measurable subset of $[0,1]$ &
$f^{-1}(\mathcal N_\rho)$, the central region in vertex space;
$\lambda(\Omega_\rho^c)=\nu(\mathcal T_\rho)$.
It is distinct from the factor-value region $\mathcal N_\rho$. \\
\end{notationtable}

\subsubsection*{\Cref{sec:concluding-remarks}: Conjectural optimizer structure}

The following symbols belong to \Cref{conj:symmetry-breaking-optimizers}.

\begin{notationtable}
$\mathcal B_H$ & A subset of $(0,1)^2$ &
$\{(p,r):p<\pc(r)\}$; the symmetry-breaking region. \\
$D_H,\pi_H$ & A nonempty interval $D_H\subseteq(0,1)$ &
$\pi_H:D_H\to(0,1)$ is the conjectured strictly increasing analytic
function parametrizing nonuniqueness by the target density $r$. \\
$\mathcal C_H$ & A curve in $\mathcal B_H$ &
$\{(\pi_H(r),r):r\in D_H\}$; the conjectured set of parameters with
nonunique optimizers, not the rank-one curve $(p_h,r_h)$. \\
$\alpha_*,\eta_{\alpha_*}$ & $0<\alpha_*\le1/2$ &
The prescribed limiting smaller-block size and a conjectured continuous
path $\eta_{\alpha_*}:(0,1)\to\mathcal B_H\setminus\mathcal C_H$
approaching $P_*$ with this limit. \\
\end{notationtable}
